% Options for packages loaded elsewhere
\PassOptionsToPackage{unicode}{hyperref}
\PassOptionsToPackage{hyphens}{url}
\PassOptionsToPackage{dvipsnames,svgnames,x11names}{xcolor}
%
\documentclass[
  10pt,
  a4paper,
]{article}
\usepackage{amsmath,amssymb}
\usepackage{iftex}
\ifPDFTeX
  \usepackage[T1]{fontenc}
  \usepackage[utf8]{inputenc}
  \usepackage{textcomp} % provide euro and other symbols
\else % if luatex or xetex
  \usepackage{unicode-math} % this also loads fontspec
  \defaultfontfeatures{Scale=MatchLowercase}
  \defaultfontfeatures[\rmfamily]{Ligatures=TeX,Scale=1}
\fi
\usepackage{lmodern}
\ifPDFTeX\else
  % xetex/luatex font selection
\fi
% Use upquote if available, for straight quotes in verbatim environments
\IfFileExists{upquote.sty}{\usepackage{upquote}}{}
\IfFileExists{microtype.sty}{% use microtype if available
  \usepackage[]{microtype}
  \UseMicrotypeSet[protrusion]{basicmath} % disable protrusion for tt fonts
}{}
\makeatletter
\@ifundefined{KOMAClassName}{% if non-KOMA class
  \IfFileExists{parskip.sty}{%
    \usepackage{parskip}
  }{% else
    \setlength{\parindent}{0pt}
    \setlength{\parskip}{6pt plus 2pt minus 1pt}}
}{% if KOMA class
  \KOMAoptions{parskip=half}}
\makeatother
\usepackage{xcolor}
\usepackage[margin=2.2cm]{geometry}
\usepackage{color}
\usepackage{fancyvrb}
\newcommand{\VerbBar}{|}
\newcommand{\VERB}{\Verb[commandchars=\\\{\}]}
\DefineVerbatimEnvironment{Highlighting}{Verbatim}{commandchars=\\\{\}}
% Add ',fontsize=\small' for more characters per line
\newenvironment{Shaded}{}{}
\newcommand{\AlertTok}[1]{\textcolor[rgb]{1.00,0.00,0.00}{\textbf{#1}}}
\newcommand{\AnnotationTok}[1]{\textcolor[rgb]{0.38,0.63,0.69}{\textbf{\textit{#1}}}}
\newcommand{\AttributeTok}[1]{\textcolor[rgb]{0.49,0.56,0.16}{#1}}
\newcommand{\BaseNTok}[1]{\textcolor[rgb]{0.25,0.63,0.44}{#1}}
\newcommand{\BuiltInTok}[1]{\textcolor[rgb]{0.00,0.50,0.00}{#1}}
\newcommand{\CharTok}[1]{\textcolor[rgb]{0.25,0.44,0.63}{#1}}
\newcommand{\CommentTok}[1]{\textcolor[rgb]{0.38,0.63,0.69}{\textit{#1}}}
\newcommand{\CommentVarTok}[1]{\textcolor[rgb]{0.38,0.63,0.69}{\textbf{\textit{#1}}}}
\newcommand{\ConstantTok}[1]{\textcolor[rgb]{0.53,0.00,0.00}{#1}}
\newcommand{\ControlFlowTok}[1]{\textcolor[rgb]{0.00,0.44,0.13}{\textbf{#1}}}
\newcommand{\DataTypeTok}[1]{\textcolor[rgb]{0.56,0.13,0.00}{#1}}
\newcommand{\DecValTok}[1]{\textcolor[rgb]{0.25,0.63,0.44}{#1}}
\newcommand{\DocumentationTok}[1]{\textcolor[rgb]{0.73,0.13,0.13}{\textit{#1}}}
\newcommand{\ErrorTok}[1]{\textcolor[rgb]{1.00,0.00,0.00}{\textbf{#1}}}
\newcommand{\ExtensionTok}[1]{#1}
\newcommand{\FloatTok}[1]{\textcolor[rgb]{0.25,0.63,0.44}{#1}}
\newcommand{\FunctionTok}[1]{\textcolor[rgb]{0.02,0.16,0.49}{#1}}
\newcommand{\ImportTok}[1]{\textcolor[rgb]{0.00,0.50,0.00}{\textbf{#1}}}
\newcommand{\InformationTok}[1]{\textcolor[rgb]{0.38,0.63,0.69}{\textbf{\textit{#1}}}}
\newcommand{\KeywordTok}[1]{\textcolor[rgb]{0.00,0.44,0.13}{\textbf{#1}}}
\newcommand{\NormalTok}[1]{#1}
\newcommand{\OperatorTok}[1]{\textcolor[rgb]{0.40,0.40,0.40}{#1}}
\newcommand{\OtherTok}[1]{\textcolor[rgb]{0.00,0.44,0.13}{#1}}
\newcommand{\PreprocessorTok}[1]{\textcolor[rgb]{0.74,0.48,0.00}{#1}}
\newcommand{\RegionMarkerTok}[1]{#1}
\newcommand{\SpecialCharTok}[1]{\textcolor[rgb]{0.25,0.44,0.63}{#1}}
\newcommand{\SpecialStringTok}[1]{\textcolor[rgb]{0.73,0.40,0.53}{#1}}
\newcommand{\StringTok}[1]{\textcolor[rgb]{0.25,0.44,0.63}{#1}}
\newcommand{\VariableTok}[1]{\textcolor[rgb]{0.10,0.09,0.49}{#1}}
\newcommand{\VerbatimStringTok}[1]{\textcolor[rgb]{0.25,0.44,0.63}{#1}}
\newcommand{\WarningTok}[1]{\textcolor[rgb]{0.38,0.63,0.69}{\textbf{\textit{#1}}}}
\setlength{\emergencystretch}{3em} % prevent overfull lines
\providecommand{\tightlist}{%
  \setlength{\itemsep}{0pt}\setlength{\parskip}{0pt}}
\setcounter{secnumdepth}{-\maxdimen} % remove section numbering
\usepackage{microtype}
\usepackage{longtable,booktabs,array}
\usepackage{xcolor}
\usepackage{fancyhdr}
\usepackage{fvextra}
\DefineVerbatimEnvironment{Highlighting}{Verbatim}{breaklines,breakanywhere,commandchars=\\\{\}}
\pagestyle{fancy}
\fancyhf{}
\fancyhead[L]{DDTA Research Methodology}
\fancyhead[R]{BA5 / Controlled Authoring}
\fancyfoot[C]{\thepage}
\setlength{\headheight}{14pt}
\emergencystretch=3em
\ifLuaTeX
  \usepackage{selnolig}  % disable illegal ligatures
\fi
\usepackage{bookmark}
\IfFileExists{xurl.sty}{\usepackage{xurl}}{} % add URL line breaks if available
\urlstyle{same}
\hypersetup{
  pdftitle={DDTA research work plan after documentation-layer closure - R22},
  colorlinks=true,
  linkcolor={blue},
  filecolor={Maroon},
  citecolor={Blue},
  urlcolor={blue},
  pdfcreator={LaTeX via pandoc}}

\title{DDTA research work plan after documentation-layer closure - R22}
\author{}
\date{}

\begin{document}
\maketitle

\section{DDTA research work plan after documentation-layer
closure}\label{ddta-research-work-plan-after-documentation-layer-closure}

\textbf{WORK PLAN - REVISION 22}

\textbf{Status:} active plan after BA5 canonical
registry/controlled-authoring closure.

\textbf{Supersedes:} Revision 21 only for forward execution state;
R1-R21 remain historical research records.

\subsection{1. Fixed prior state}\label{1-fixed-prior-state}

\begin{itemize}
\tightlist
\item
  Chapters 2-4: CLOSED / FINAL for current thesis scope.
\item
  Documentation layer: CLOSED.
\item
  BA0 responsibility/non-goals: CLOSED.
\item
  BA1 minimal BAE identity ontology: CLOSED.
\item
  BA2 relation/action vocabulary: CLOSED BY BA2-T4.
\item
  BA3 provenance/derivation/identity/lifecycle/change-revalidation:
  CLOSED BY BA3-T4.
\item
  BA4 projection boundary/traceability/coverage/interpretation: CLOSED
  BY BA4-T3.
\item
  BA5 canonical semantic registry/controlled authoring: CLOSED BY
  BA5-T3.
\item
  \texttt{BAReferent} and \texttt{BAProposition}: ACCEPTED first-class
  semantic identity families.
\item
  ThreatForge remains a case study/tooling instrument, not DDTA semantic
  authority.
\end{itemize}

\subsection{2. BA5 closure carried
forward}\label{2-ba5-closure-carried-forward}

BA5 closes the strict controlled-authoring hypothesis:

\begin{Shaded}
\begin{Highlighting}[]
\NormalTok{semantically operative DDTA binding}
\NormalTok{    {-}\textgreater{} exact canonical token/name}

\NormalTok{unregistered alias / synonym}
\NormalTok{    {-}\textgreater{} not normative semantic input}

\NormalTok{free explanatory prose}
\NormalTok{    {-}\textgreater{} allowed outside semantic bindings}
\end{Highlighting}
\end{Shaded}

The closed registry responsibilities are domain-scoped: referent names,
BA2 operators, operator-scoped roles, semantic kinds and
controlled-value domains remain distinct in authority/scope even when
physically co-located.

Normative synonym machinery is \textbf{REJECTED AS UNNECESSARY} for the
current thesis scope. Optional lexical search, translation, migration or
NLP/LLM assistance remains \textbf{DEFERRED} and must not become
semantic authority.

\subsection{3. BA5 extension/reopen
rule}\label{3-ba5-extensionreopen-rule}

Do not reopen BA5 for convenience.

Reopen only the smallest BA5 responsibility on concrete evidence that
exact canonical bindings plus presentation-local wording cannot preserve
required shared meaning or reproducibility.

A materially new operator or role contract reopens BA2, not BA5. A new
independently reusable identity applies the BA1 split criterion.
Method-specific terminology remains downstream.

\subsection{4. BA6 objective - complete Base Analysis integrated
regression and
closure}\label{4-ba6-objective---complete-base-analysis-integrated-regression-and-closure}

\textbf{Status: NOT STARTED / NEXT PHASE.}

BA6 is the integration/adversarial closure phase for the complete
BA0-BA5 design. It is not a new ontology-building phase by default.

The integrated chain under test is:

\begin{Shaded}
\begin{Highlighting}[]
\NormalTok{governed documentation}
\NormalTok{ {-}\textgreater{} canonical registry{-}controlled authoring}
\NormalTok{ {-}\textgreater{} accepted Base Analysis}
\NormalTok{ {-}\textgreater{} BA3 provenance / review / continuity}
\NormalTok{ {-}\textgreater{} multiple BA4 projections}
\NormalTok{ {-}\textgreater{} governed project change}
\NormalTok{ {-}\textgreater{} BA revalidation / rebuild}
\NormalTok{ {-}\textgreater{} projection rebuild / comparison}
\end{Highlighting}
\end{Shaded}

BA6 must determine whether the closed responsibilities work together
without requiring consumers to reconstruct missing shared
methodology-neutral meaning from raw prose.

\subsection{5. BA6 evidence minimum}\label{5-ba6-evidence-minimum}

The phase must include:

\begin{enumerate}
\def\labelenumi{\arabic{enumi}.}
\tightlist
\item
  regression over the closed facial-access corpus and mutation
  pressures;
\item
  regression over the closed order/WMS/provider corpus and change
  pressures;
\item
  at least one structurally different holdout not used to tune the
  BA0-BA5 lower bounds;
\item
  canonical-name and registry-revision resolution through the complete
  chain;
\item
  provenance/source drill-down from projection to BA to governed source;
\item
  change localization, BA3 continuity/revalidation and projection
  rebuild;
\item
  at least two materially different consumers/projections; and
\item
  explicit search for a material counterexample requiring the smallest
  BA0-BA5 reopen.
\end{enumerate}

The holdout is for integration falsification, not for silently expanding
DDTA into a universal systems-modeling language.

\subsection{6. BA6 guardrails}\label{6-ba6-guardrails}

Do not:

\begin{itemize}
\tightlist
\item
  add a new BAE family without a material BA1 counterexample;
\item
  add a BA2 operator/role merely for convenience;
\item
  weaken BA3 source/provenance/change accountability;
\item
  let a projection repair missing shared BA meaning by rereading raw
  prose;
\item
  let method taxonomy redefine BA semantics;
\item
  reintroduce synonym normalization without a BA5 material
  counterexample;
\item
  treat ThreatForge runtime classes as DDTA semantics;
\item
  start AnalysisRecord, Common Finding or formal STRIDE/STRIDE-AI schema
  design; or
\item
  declare Base Analysis closed before the integrated regression passes.
\end{itemize}

\subsection{7. BA6 counterexample
routing}\label{7-ba6-counterexample-routing}

Use the smallest-responsibility rule:

\begin{Shaded}
\begin{Highlighting}[]
\NormalTok{material counterexample}
\NormalTok{    {-}\textgreater{} identify exact failed responsibility}
\NormalTok{    {-}\textgreater{} reopen only that BA phase/contract}
\NormalTok{    {-}\textgreater{} repair and rerun affected regression}
\end{Highlighting}
\end{Shaded}

Examples:

\begin{itemize}
\tightlist
\item
  missing independent semantic identity -\textgreater{} BA1;
\item
  missing method-neutral proposition/operator/role semantics
  -\textgreater{} BA2;
\item
  unreproducible provenance/change continuity -\textgreater{} BA3;
\item
  projection coverage/trace/interpretation failure -\textgreater{} BA4;
\item
  canonical registry/authoring failure -\textgreater{} BA5.
\end{itemize}

Do not redesign Base Analysis wholesale from one local failure.

\subsection{8. BA6 exit condition}\label{8-ba6-exit-condition}

BA6 may close the complete Base Analysis milestone only if:

\begin{enumerate}
\def\labelenumi{\arabic{enumi}.}
\tightlist
\item
  the closed corpora regress without material semantic loss;
\item
  the structurally different holdout does not force an unresolved
  shared-semantic gap;
\item
  source drill-down and historical reproducibility remain intact;
\item
  governed change can be localized and revalidated without whole-BA
  invalidation by default;
\item
  multiple projections can be rebuilt without becoming project
  authority;
\item
  canonical registry/controlled-authoring semantics survive integration;
  and
\item
  any discovered counterexample has been routed to and resolved in the
  smallest responsible phase.
\end{enumerate}

Only then may the state become:

\begin{Shaded}
\begin{Highlighting}[]
\NormalTok{BASE ANALYSIS    CLOSED FOR CURRENT THESIS SCOPE}
\end{Highlighting}
\end{Shaded}

\subsection{9. Downstream analysis envelope remains
blocked}\label{9-downstream-analysis-envelope-remains-blocked}

Until BA6 closes, do not start:

\begin{itemize}
\tightlist
\item
  A1 - AnalysisRecord / execution envelope;
\item
  A2 - common reviewed Finding boundary;
\item
  A3 - analysis change/provenance integration;
\item
  formal STRIDE/STRIDE-AI overlay schemas; or
\item
  ThreatForge implementation work that would imply unclosed semantic
  authority.
\end{itemize}

\subsection{10. Next authorized work}\label{10-next-authorized-work}

The next authorized phase is:

\begin{quote}
\textbf{BA6 - complete Base Analysis integrated regression and closure.}
\end{quote}

Before execution, define the smallest first BA6 microstep against the
closed BA0-BA5 contract. Revision 22 intentionally does not invent a BA6
microstep count or execute BA6 inside the BA5-T3 closure package.

\end{document}
